\chapter{Introduction}\label{c:introduction}

Recovering the three-dimensional structure of a scene from visual observations
is a foundational problem in computer vision, with applications spanning robotics,
autonomous navigation, and augmented reality~\cite{hartley2004multiple}. In
particular, metric depth estimation enables downstream systems to reason about
geometry and spatial relationships in real environments. While deep learning has
substantially improved monocular depth estimation~\cite{eigen2014depth,
bhat2023zoedepth}, extending these methods to sparse multi-view settings remains
a challenging and practically important problem.

When multiple synchronized cameras observe the same scene over time, the
reconstruction problem extends naturally from static 3D geometry to modeling
temporal scene evolution - commonly referred to as 4D reconstruction. Here the
objective is not only to recover accurate per-frame geometry, but to ensure that
the recovered structure is consistent across the sequence. Achieving both
geometric accuracy and temporal stability under sparse camera configurations
(typically two to four cameras, as encountered in robotic platforms) is the
central problem of this thesis.

\section{Background and Motivation}

Classical approaches to multi-view reconstruction, such as Structure-from-Motion
(SfM)~\cite{schoenberger2016sfm} and Multi-View Stereo
(MVS)~\cite{schoenberger2016mvs}, rely on sparse feature correspondences and
struggle in scenes with textureless regions, dynamic objects, or limited
viewpoint overlap. Recent learning-based methods have largely superseded these
pipelines by reformulating reconstruction as dense prediction.
DUSt3R~\cite{dust3r} and MASt3R~\cite{mast3r} in particular demonstrate strong
robustness in sparse-view settings by directly regressing pointmaps from image
pairs, bypassing explicit camera calibration. MASt3R-SfM~\cite{mast3rsfm}
extends this framework to full scene reconstruction through global pointmap alignment, whereas VGGT~\cite{vggt} employs a joint optimization of camera parameters and dense geometry across multi-view collections to ensure global consistency.

\todo[inline]{Needs citation}These spatial methods are typically evaluated on a per-frame basis, and their
behavior in the temporal domain is less well understood. Methods designed
specifically for dynamic scenes - such as VGGT4D~\cite{vggt4d} and
Mono4D~\cite{mono4d} - incorporate temporal information to encourage smoother
reconstructions, but may trade off spatial precision in doing so. Whether methods
that excel under standard per-frame benchmarks also produce temporally stable 4D
reconstructions remains an open and underexplored question.

\section{Problem Statement and Contributions}

This thesis investigates metric 4D reconstruction from sparse multi-view video.
Given synchronized video streams from two to four cameras, the goal is to produce
scene reconstructions that are both geometrically accurate and temporally stable.
A particular focus is placed on characterizing the trade-off between these two
desiderata - a gap that is obscured by conventional per-frame evaluation
protocols - and on understanding which design choices drive or reduce temporal
inconsistency.

The main contributions of this work are:

\begin{itemize}
    \item \textbf{Evaluation framework.} A benchmarking pipeline for sparse
    multi-view video depth estimation that measures both per-frame spatial
    accuracy and sequence-level temporal consistency across multiple datasets
    \todo[inline]{Datasets require citation. Datasets subject to change}
    (DexYCB, Hi4D, 4D-Dress, MultiCamRobolab).

    \item \textbf{Consistency metric.} A quantitative measure of the spatial--
    temporal error gap,
    \[
      \Delta_{\text{consistency}} = \text{Chamfer}_{4D} - \text{Chamfer}_{3D},
    \]
    which captures how much per-frame alignment masks underlying temporal
    instability (formally introduced in \Cref{c:methodology}).

    \item \textbf{Empirical analysis.} A systematic comparison of spatial methods
    (MASt3R, VGGT) and temporal methods (VGGT4D, Mono4D) across two to
    four camera configurations, exposing failure modes and trade-offs not visible
    under standard benchmarks.

    \item \textbf{Alignment and refinement study.} An investigation of temporal
    alignment strategies and a global pose-graph refinement approach aimed at
    improving reconstruction stability without sacrificing spatial fidelity.
\end{itemize}

\section{Thesis Outline}

\Cref{c:related_work} reviews prior work on depth estimation, multi-view
reconstruction, and dynamic scene modeling. \Cref{c:methodology} describes the
datasets, metrics, and alignment strategies used in this study.
\Cref{c:experiments} presents quantitative and qualitative results.
\Cref{c:limitations} analyzes the key findings and limitations, and
\Cref{c:conclusions} summarizes the contributions and outlines directions for
future work.